\section{Background}

The emergence of multi-agent systems as a dominant paradigm for orchestrating complex, autonomous workloads necessitates rigorous frameworks for agent lifecycle management, accountability, and provenance. This section situates the BX3 Framework within the broader landscape of recursive spawning architectures, hierarchical task decomposition, and immutable delegation primitives. We survey prior art across five axes: agent lifecycle management, accountability and provenance, delegation chain semantics, hierarchical decomposition, and substrate requirements for immutable parent chains.

\subsection{Agent Lifecycle Management in Multi-Agent Systems}

Agent lifecycle management encompasses the creation, provisioning, coordination, and decommissioning of autonomous agents. Prior work establishes that lifecycle management in distributed AI systems is fundamentally distinct from traditional microservice orchestration due to agents' autonomous goal-directed behavior and emergent sub-task generation~\cite{aitlp-draft}. The Agent Identity, Trust and Lifecycle Protocol (AITLP) defines formal lifecycle states---spawned, active, suspended, revoked---and hierarchical mandate enforcement structures that govern what agents are permitted to do across generational boundaries~\cite{aitlp-draft}. This provides a principled state machine for agent existence from instantiation through termination.

Dynamic spawning architectures extend this model with runtime-aware spawning policies. AgentSpawn demonstrates that adaptive, complexity-driven spawning---where agents are instantiated on-demand based on measured task difficulty---yields approximately 34\% improvement in completion rates on long-horizon benchmarks compared to static pre-allocation~\cite{agentspawn}. Crucially, AgentSpawn incorporates memory transfer mechanisms enabling child agents to inherit relevant context at spawn time, addressing the continuity problem in generational agent handoffs~\cite{agentspawn}. AGENTHIVE advances this by treating delegation as a first-class primitive: any agent may spawn and coordinate sub-agents, enabling decentralized control without a fixed central orchestrator and allowing topology to emerge dynamically from task demands~\cite{agenthive}.

ReDel provides an open-source toolkit for recursive multi-agent systems with event-driven logging, replay, and visualization of delegation graphs, supporting fine-grained inspection of agent state transitions across the full lifecycle~\cite{redel}. Together, these systems establish that lifecycle management in multi-agent environments requires (1) formal state models governing spawn and revocation, (2) runtime-adaptive spawning policies, (3) memory and context transfer at agent creation, and (4) event-level logging for post-hoc reconstruction.

\subsection{Accountability and Provenance in Distributed Systems}

As agents proliferate across distributed pipelines, accountability and provenance become critical infrastructure requirements. The fundamental distinction is between accountability (who is responsible and compliant) and auditability (the properties that enable reliable auditing)~\cite{auditable-agents}. Auditable Agents formalizes five dimensions of agent auditability: action recoverability, lifecycle coverage, policy checkability, responsibility attribution, and evidence integrity, arguing that a single mechanism class (detect, enforce, or recover) is insufficient in isolation~\cite{auditable-agents}. This framework establishes that accountability infrastructure must span the full agent lifecycle with tamper-evident records and layered evidence.

Several architectural approaches address this at the protocol level. PROV-AGENT extends the W3C PROV model with agent-centric metadata capturing prompts, responses, and decisions, linking them to end-to-end workflow provenance across federated environments~\cite{prov-agent}. The Human Delegation Provenance (HDP) protocol implements cryptographic, offline-verifiable provenance by binding human authorization to AI orchestration sessions via an append-only hop chain, where each delegation step is signed and verifiable without network lookups or third-party trust anchors~\cite{hdp-protocol}. The Warrant Certificate Authority (WCA) framework addresses epistemic provenance in tool-call chains by certifying data sources rather than data content, analogous to PKI-style trust models, and proposes Warrant Attestation Levels (WAL-0 to WAL-3) as a staged adoption path for provenance guarantees~\cite{wca}.

TrustTrack embeds cryptographic identity, policy commitments, and tamper-resistant behavioral logs directly into agent infrastructure via a four-layer Trust-Native Stack, reframing compliance as a design constraint rather than an afterthought~\cite{trusttrack}. Similarly, the Verifiable AI Provenance (VAP) framework composes existing IETF primitives (SCITT for transparency, RATS for remote attestation, COSE for signing) into an evidentiary-grade decision trail architecture without introducing new cryptographic primitives~\cite{vap-framework}. Auditable tool-call logging via Ed25519-signed hash chains provides a lightweight, offline-verifiable method for establishing tamper-evident provenance in agentic tool interactions~\cite{langchain-audit}.

Taken together, the field establishes that provenance in multi-agent systems must be cryptographic, tamper-evident, and span generational boundaries. The central tension is between the practical overhead of cryptographic chaining and the requirement for low-latency, high-throughput agent execution.

\subsection{Fixed vs. Mutable Delegation Chains}

A delegation chain records the sequence of principal-to-agent (or agent-to-agent) authorizations enabling a given action. The fundamental architectural choice is whether this chain is fixed at spawn time or mutable during execution.

Fixed delegation chains offer strong provenance guarantees: because each hop is immutably recorded, the full chain of custody is traceable after the fact. However, fixed chains are brittle in dynamic environments where task requirements shift mid-execution, and they cannot accommodate adaptive re-decomposition when a spawned agent encounters unexpected subtask complexity. The HDP protocol and WCA framework both assume append-only semantics precisely because mutability after the fact destroys evidentiary value.

Mutable delegation chains, as implemented in AGENTHIVE and AgentSpawn, allow agents to re-delegate, escalate, or cascade authority at runtime based on task demands~\cite{agenthive,agentspawn}. This flexibility improves solution space coverage but introduces a fundamental accountability gap: if a delegation chain can be rewritten after the fact, provenance becomes unreliable. ADAPT's as-needed decomposition further complicates this by enabling agents to recursively spawn sub-tasks without pre-specified depth limits, making the delegation topology itself an emergent rather than a planned artifact~\cite{adapt}.

The Recursive Delegation Engine (RDE) addresses this by representing tasks as a delegation topology (directed graph) where leaf nodes execute and propagate binary rewards (success/failure) back up the chain, informing future delegation decisions~\cite{rde}. RDE's bandit-based selection policies optimize delegation dynamically but do not inherently preserve an immutable record of the topology that produced a given outcome. This tension between adaptive mutability and immutable provenance is unresolved in the literature and motivates the BX3 Framework's separation of execution topology from immutable parent-chain substrate.

\subsection{Hierarchical Task Decomposition in AI}

Hierarchical Task Decomposition (HTD) is foundational to scalable multi-agent planning. Formal models include Hierarchical Task Networks (HTNs), which decompose goals into methods and operators organized as trees or DAGs with ordering and resource constraints, and Hierarchical MDPs / Semi-Markov Decision Processes, which represent root tasks as high-level processes invoking subtasks with their own state, action, and reward structures~\cite{htd-survey}. These models enable value-function reuse across decomposition levels (e.g., MAXQ, options) and support transfer of knowledge across tasks and environments~\cite{htd-survey}.

ADAPT (As-Needed Decomposition and Planning with Language Models) demonstrates that recursive, on-demand decomposition yields approximately 33\% improvement over prior state-of-the-art by avoiding both over- and under-decomposition: agents expand sub-plans only when a subtask proves challenging, balancing planning granularity against computational cost~\cite{adapt}. AgentOrchestra instantiates this as a central planning agent that decomposes objectives into sub-goals and delegates to specialized, tool-enabled sub-agents, empirically demonstrating success rate gains over flat-agent baselines on real-world benchmarks~\cite{agentorchestra}. The sub-agent receives minimal context to reduce token bloat and hallucination risk~\cite{agentorchestra}, a pattern echoed in SubAgent architectures where parent agents retain only high-signal summaries from spawned agents rather than step-by-step traces.

Critically, hierarchical decomposition creates a structural duality: the decomposition tree (task structure) and the delegation topology (agent structure) are distinct but coupled. Prior work conflates them, treating decomposition as inseparable from agent spawning. The BX3 Framework decouples these concerns by maintaining an immutable parent chain substrate independent of task decomposition structure, enabling independent verification of both what was done (decomposition) and who authorized it (delegation).

\subsection{The BX3 Framework as Substrate for Immutable Parent Chains}

The foregoing survey reveals an unresolved gap: no existing framework simultaneously achieves (1) formal agent lifecycle management with state-machine governance, (2) cryptographic, immutable provenance across generational boundaries, (3) adaptive delegation topology without mutability destroying auditability, and (4) decoupled decomposition and delegation structures.

The BX3 (Brodi Blanco Bounded Bidirectional Block) Framework addresses this gap by establishing an immutable parent-chain substrate as the foundational layer for agent spawning. Every agent in a BX3-mediated system is spawned with a cryptographically signed reference to its parent agent, creating an append-only, tamper-evident chain of custody that is independent of task decomposition structure. The parent chain records only authorization provenance---which agent spawned which, under what mandate, at what timestamp---rather than task content, enabling verifiable accountability without constraining decomposition flexibility.

The substrate layer provides three guarantees: First, parent-chain immutability is enforced cryptographically at the substrate level, making the chain of custody tamper-evident regardless of runtime mutation in agent behavior or delegation topology. Second, the substrate decouples provenance recording from task execution, allowing agents to engage in adaptive, recursive decomposition (ADAPT-style, HTN-style, or hybrid) without modifying the authorization record. Third, the BX3 Framework integrates with formal lifecycle state models (per AITLP) and provenance logging (per PROV-AGENT and HDP), positioning the parent chain as the cryptographic spine connecting lifecycle state transitions to evidentiary audit trails.

This separation of concerns---authorization provenance in the substrate, task decomposition in agent execution---enables the BX3 Framework to support verifiable accountability in recursive spawning scenarios where prior approaches sacrifice either adaptability or auditability. We examine this claim formally in Section~\ref{sec:framework} and evaluate it empirically in Section~\ref{sec:evaluation}.

\bibliographystyle{plain}
\bibliography{references}
